% GENERATED FILE. Edit canonical lessons, then run npm run generate:books.
% canonical-source-hash: fnv1a64:a0640f0a8d85d365
% canonical-lessons: IT-C464-cinghiale, IT-C464-cervo, IT-C464-lepre, IT-C464-falco, IT-C464-brucare

\chapter{Animali selvatici}
\label{ch:it-animali-selvatici}
\hlchaptermodality{464}

\begin{chapteropening}
\textbf{By the end of this chapter:} \emph{I can name wild animals I might see: the wild boar, the deer, the hare, the falcon, and animals grazing.}

Complete the last lesson of chapter 464: I can name wild animals I might see: the wild boar, the deer, the hare, the falcon, and animals grazing.
\end{chapteropening}

\section[il cinghiale]{il cinghiale — the wild boar}

\label{lesson:IT-C464-cinghiale}



\begin{quote}
Before the new one: say the Italian for the tide, then the Italian for to dive.
\end{quote}

\subsection*{You'll want to know: il cinghiale}
\textbf{il cinghiale} — \textquotedblleft{}the wild boar\textquotedblright{}.

\begin{quote}
Di notte i cinghiali scendono in paese.
\end{quote}

\begin{quote}
\emph{At night the boars come down into the village.}
\end{quote}

Common in the hills, and a popular sauce for pasta.

\begin{cousinweb}
\emph{cinghiale} is usually traced to Late Latin \emph{singulāris}, short for \emph{porcus singulāris} 'the pig that lives alone', from \emph{singulāris} 'alone, solitary' (Lewis \& Short); it is inherited, worn down in speech.
\end{cousinweb}

\begin{grammarlens}[title={What you've built}]
One more word for telling, explaining and describing at B1.
\end{grammarlens}

\subsection*{Your turn}
\begin{itemize}
\raggedright
  \item \emph{Say it:} \emph{il cinghiale}
  \item \emph{Say it:} \emph{il cinghiale}, once more
  \item \emph{Say it:} a sentence of your own with \emph{il cinghiale}
  \item \emph{Recall:} say the Italian for the tide, then the Italian for to dive, then say \emph{il cinghiale} again
\end{itemize}

\begin{tcolorbox}[breakable,colback=teal!4,colframe=teal!35!black,title={Before you move on}]
What does il cinghiale mean? (\textquotedblleft{}The wild boar\textquotedblright{}.) Say it once more.
\end{tcolorbox}
\section[il cervo]{il cervo — the deer, the stag}

\label{lesson:IT-C464-cervo}



\begin{quote}
Before the new one: say the Italian for to dive, then the Italian for the wild boar.
\end{quote}

\subsection*{You'll want to know: il cervo}
\textbf{il cervo} — \textquotedblleft{}the deer, the stag\textquotedblright{}.

\begin{quote}
Nel parco abbiamo visto un cervo.
\end{quote}

\begin{quote}
\emph{In the park we saw a deer.}
\end{quote}

\emph{Il capriolo} is the smaller roe deer.

\begin{cousinweb}
\emph{cervo} is inherited from Latin \emph{cervus} 'deer' (EtymDB).
\end{cousinweb}

\begin{grammarlens}[title={What you've built}]
One more word for telling, explaining and describing at B1.
\end{grammarlens}

\subsection*{Your turn}
\begin{itemize}
\raggedright
  \item \emph{Say it:} \emph{il cervo}
  \item \emph{Say it:} \emph{il cervo}, once more
  \item \emph{Say it:} a sentence of your own with \emph{il cervo}
  \item \emph{Recall:} say the Italian for to dive, then the Italian for the wild boar, then say \emph{il cervo} again
\end{itemize}

\begin{tcolorbox}[breakable,colback=teal!4,colframe=teal!35!black,title={Before you move on}]
What does il cervo mean? (\textquotedblleft{}The deer, the stag\textquotedblright{}.) Say it once more.
\end{tcolorbox}
\section[la lepre]{la lepre — the hare}

\label{lesson:IT-C464-lepre}



\begin{quote}
Before the new one: say the Italian for the wild boar, then the Italian for the deer.
\end{quote}

\subsection*{You'll want to know: la lepre}
\textbf{la lepre} — \textquotedblleft{}the hare\textquotedblright{}.

\begin{quote}
Una lepre attraversa il campo.
\end{quote}

\begin{quote}
\emph{A hare is crossing the field.}
\end{quote}

Bigger than \emph{il coniglio}, the rabbit.

\begin{cousinweb}
\emph{lepre} is inherited from Latin \emph{lepus, leporem} 'hare' (EtymDB).
\end{cousinweb}

\begin{grammarlens}[title={What you've built}]
One more word for telling, explaining and describing at B1.
\end{grammarlens}

\subsection*{Your turn}
\begin{itemize}
\raggedright
  \item \emph{Say it:} \emph{la lepre}
  \item \emph{Say it:} \emph{la lepre}, once more
  \item \emph{Say it:} a sentence of your own with \emph{la lepre}
  \item \emph{Recall:} say the Italian for the wild boar, then the Italian for the deer, then say \emph{la lepre} again
\end{itemize}

\begin{tcolorbox}[breakable,colback=teal!4,colframe=teal!35!black,title={Before you move on}]
What does la lepre mean? (\textquotedblleft{}The hare\textquotedblright{}.) Say it once more.
\end{tcolorbox}
\section[il falco]{il falco — the falcon, the hawk}

\label{lesson:IT-C464-falco}



\begin{quote}
Before the new one: say the Italian for the deer, then the Italian for the hare.
\end{quote}

\subsection*{You'll want to know: il falco}
\textbf{il falco} — \textquotedblleft{}the falcon, the hawk\textquotedblright{}.

\begin{quote}
Il falco vola sopra la valle.
\end{quote}

\begin{quote}
\emph{The falcon flies over the valley.}
\end{quote}

\emph{Occhio di falco} — eagle eye.

\begin{cousinweb}
\emph{falco} is inherited from Late Latin \emph{falcō} 'falcon' (EtymDB).
\end{cousinweb}

\begin{grammarlens}[title={What you've built}]
One more word for telling, explaining and describing at B1.
\end{grammarlens}

\subsection*{Your turn}
\begin{itemize}
\raggedright
  \item \emph{Say it:} \emph{il falco}
  \item \emph{Say it:} \emph{il falco}, once more
  \item \emph{Say it:} a sentence of your own with \emph{il falco}
  \item \emph{Recall:} say the Italian for the deer, then the Italian for the hare, then say \emph{il falco} again
\end{itemize}

\begin{tcolorbox}[breakable,colback=teal!4,colframe=teal!35!black,title={Before you move on}]
What does il falco mean? (\textquotedblleft{}The falcon, the hawk\textquotedblright{}.) Say it once more.
\end{tcolorbox}
\section[brucare]{brucare — to graze, to nibble (grass or leaves)}

\label{lesson:IT-C464-brucare}



\begin{quote}
Before the new one: say the Italian for the hare, then the Italian for the falcon.
\end{quote}

\subsection*{You'll want to know: brucare}
\textbf{brucare} — \textquotedblleft{}to graze, to nibble (grass or leaves)\textquotedblright{}.

\begin{quote}
Le capre brucano l'erba.
\end{quote}

\begin{quote}
\emph{The goats graze on the grass.}
\end{quote}

Of animals eating leaves and grass.

\begin{cousinweb}
The origin of \emph{brucare} is not settled; it may be linked to an old word for a shoot or a twig.
\end{cousinweb}

\begin{grammarlens}[title={What you've built}]
One more word for telling, explaining and describing at B1.
\end{grammarlens}

\subsection*{Your turn}
\begin{itemize}
\raggedright
  \item \emph{Say it:} \emph{brucare}
  \item \emph{Say it:} \emph{brucare}, once more
  \item \emph{Say it:} a sentence of your own with \emph{brucare}
  \item \emph{Recall:} say the Italian for the hare, then the Italian for the falcon, then say \emph{brucare} again
\end{itemize}

\begin{tcolorbox}[breakable,colback=teal!4,colframe=teal!35!black,title={Before you move on}]
What does brucare mean? (\textquotedblleft{}To graze, to nibble (grass or leaves)\textquotedblright{}.) Say it once more.
\end{tcolorbox}
